\section{Introduction}
\label{sec:introduction}

Supervised learning seeks to approximate an unknown functional relationship between a set of descriptors characterizing a system and a target property of interest. The difficulty of this approximation depends on the complexity of the underlying feature–target mapping. If the target varies smoothly over the representation space, relatively simple models can interpolate accurately and generalize well. If the mapping is highly irregular or poorly aligned with the representation, even highly expressive models may struggle to generalize beyond the training distribution.

Consequently, improving generalization has traditionally focused on designing better representations, more appropriate inductive biases, and more expressive learning architectures, with comparatively little attention given to the target itself. The target is typically regarded as a fixed quantity provided by the dataset, while preprocessing is introduced primarily to improve optimization or numerical stability. This implicitly assumes that the learning problem is determined entirely by the representation and model architecture, and that the target merely specifies what should be predicted.

Machine-learning interatomic potentials (MLIPs) provide a representative example of this broader perspective. Considerable effort has been devoted to improving representations and inductive biases in order to capture increasingly complex components of the atomistic total energy. Rich geometric descriptors, equivariant message passing, many-body interactions, and highly expressive graph neural networks have substantially improved predictive performance by making the mapping from atomic structure to energy easier to learn \cite{behler2007generalized,bartok2010gap,schutt2017schnet,gasteiger2020dimenet,gasteiger2021gemnet,batzner2022nequip,batatia2022mace}.

In contrast, comparatively little attention has been given to the target side of the learning problem. Beyond standard normalization, which merely rescales the labels, MLIP training routinely applies two additional preprocessing steps before learning \cite{schutt2023schnetpack2,pelaez2024torchmdnet2,tran2023oc22,levine2025omol25,wood2025uma,sriram2025odac25}. First, isolated atomic reference energies are subtracted from the DFT total energy to produce a reference-corrected target. Second, the remaining reference-corrected energies are adjusted using a simple linear regression based on elemental composition, often referred to as the Bag-of-Elements (BoE) correction. These steps are generally presented as preprocessing procedures intended to improve optimization, convergence, and numerical stability.

A closer examination reveals that they do considerably more. Subtracting atomic reference energies transforms the learning target from total energy to atomization (or formation-like) energy, while the subsequent BoE correction further redefines the target by removing an additional compositional contribution. These operations therefore do not simply rescale the target—they redefine the prediction task itself. The learning problem changes because the model is no longer trained on the original DFT total energy, but on a different physical quantity.

This distinction becomes particularly significant in light of our empirical observations. As shown below, applying the BoE correction substantially improves not only in-domain generalization but also out-of-distribution transferability of GemNet \cite{gasteiger2021gemnet}, despite the model already containing approximately 1.7 million trainable parameters. At first glance, this additional linear-regression step appears almost trivial compared to the expressive capacity of the graph neural network. Yet its impact on predictive performance is surprisingly large. This immediately raises a fundamental question: why can such a simple preprocessing step produce improvements that are not recovered by a highly expressive model during training?

One possible explanation is that the correction simply reduces the variance of the learning target, thereby making optimization easier. While this interpretation is plausible, it immediately raises further questions. Would any baseline that reduces target variance produce the same improvement? Is variance reduction alone sufficient to explain the observed gains in transferability? Or does the choice of baseline fundamentally redefine the learning problem in a way that better aligns the target with the model’s inductive bias?

Resolving these questions is essential for moving toward informed target-design decisions for generalizable machine learning. If target transformations do more than improve optimization, then target design is not merely a preprocessing step but a fundamental component of model design. We propose that target design should be regarded as a third design axis of supervised learning, alongside representation design and architectural design. Rather than treating the target as fixed, it should be engineered to align the learning problem with the model's inductive bias and maximize transferability. In this work, we examine this hypothesis through the lens of delta learning, seeking general principles that connect baseline choice, target landscape, and model generalization.
